\documentclass[unknownkeysallowed]{beamer}
\usepackage[utf8]{inputenc}
\usepackage{beamer_tom}
\graphicspath{{./images/}{./sharedimages/}}


\usepackage{pythonhighlight}

\usepackage{dirtree}
\newcommand{\Benchopt}{\texttt{Benchopt}}
\newcommand{\TutorialRepoURL}{https://github.com/tomMoral/26-05\_benchopt\_tuto\_icms}
\newcommand{\TutorialBenchDir}{\texttt{denoising}}
\newcommand{\BaseSolver}{\texttt{DRUNet}}
\newcommand{\NewSolver}{\texttt{DnCNN}}
\newcommand{\NewBenchmark}{\texttt{blind\_deblur}}

% -----------------------------------------------------------------------------
% RETARGETING NOTES (for another tutorial repo/benchmark)
% 1) Update these macros first:
%    - \TutorialRepoURL
%    - \TutorialBenchDir
%    - \BaseSolver / \NewSolver
%    - \NewBenchmark
% 2) Update install/run commands in the "Install and run the benchmark" frame.
% 3) Update the file paths and metric names in:
%    - "What happens during install and run?"
%    - Solver Step 3 validation checklist
%    - New benchmark Step 2 protocol changes
% 4) Update benchmark-specific solver filenames in "Step 3: add solvers and
%    validate" for the new benchmark section.
% -----------------------------------------------------------------------------

\definecolor{tomcolor}{rgb}{0.,.3,.6}
\lstset{basicstyle=\color{tomcolor}\scriptsize\ttfamily}


\author{Benchopt contributors}
\collaborators[none]{\vskip-1em}
\title{
    Hands-on \texttt{benchopt}\\
    Imaging benchmarks tutorial
}
\setbeamertemplate{title page}[frame]
\def\extraLogo{\includegraphics[height=1.5cm]{logo_benchopt.png}}


\begin{document}

%%%%%%%%%%%%%%%%%%%%%%%%%%%%%%%%%%%%%%%%%%%%%%%%%%%%%%%%%%%%%%%%%%%%%%%%%%%%%%%
%%%%%%%%%%%%%%%%%%%%%%             Headers               %%%%%%%%%%%%%%%%%%%%%%
%%%%%%%%%%%%%%%%%%%%%%%%%%%%%%%%%%%%%%%%%%%%%%%%%%%%%%%%%%%%%%%%%%%%%%%%%%%%%%%


%%%%%%%%%%%%%%%%%%%%%%%%%%%%%%%%%%%%%%%%%%%%%%%%%%%%%%%%%%%%%%%%%%%%%%%%%%%%%%%
    \begin{frame}
        \titlepage
    \end{frame}

%%%%%%%%%%%%%%%%%%%%%%%%%%%%%%%%%%%%%%%%%%%%%%%%%%%%%%%%%%%%%%%%%%%%%%%%%%%%%%%

% %%%%%%%%%%%%%%%%%%%%%%%%%%%%%%%%%%%%%%%%%%%%%%%%%%%%%%%%%%%%%%%%%%%%%%%%%%%%%%%
% \begin{frame}{Contributors from...}
%     \centering
%     \includegraphics[width=100px]{logo_inria_horiz.pdf} \hspace{1mm}
%     \includegraphics[height=45px]{um_logo.pdf} \hspace{1mm}
%     \includegraphics[height=35px]{logo_berkeley.png} \hspace{1mm}
%     \\[7mm]
%     \includegraphics[height=50px]{logo_cnrs.pdf} \hspace{5mm}
%     \includegraphics[height=50px]{logo_lund} \hspace{5mm}
%     \includegraphics[height=50px]{logo_telecom.pdf} \hspace{5mm}
%     \includegraphics[height=50px]{logo_univ_luxembourg.pdf} \\[4mm]
%     \includegraphics[height=40px]{logo_ens.png}
%     % \includegraphics[height=40px]{}
% \end{frame}
% %%%%%%%%%%%%%%%%%%%%%%%%%%%%%%%%%%%%%%%%%%%%%%%%%%%%%%%%%%%%%%%%%%%%%%%%%%%%%%%



%%%%%%%%%%%%%%%%%%%%%%%%%%%%%%%%%%%%%%%%%%%%%%%%%%%%%%%%%%%%%%%%%%%%%%%%%%%%%%%
% \begin{frame}{Why benchmarks matter}
%
%     Benchmarks are central to empirical progress in ML:\\[.5em]
%     \begin{itemize}
%         \item They define what we compare and report.
%         \item They make progress measurable across papers.
%         \item They provide the baselines used by the whole community.
%     \end{itemize}
%
%     \pause
%     \vskip2em
%     \strongpoint{No reliable benchmark, no reliable scientific claim in ML.}
%
% \end{frame}
%%%%%%%%%%%%%%%%%%%%%%%%%%%%%%%%%%%%%%%%%%%%%%%%%%%%%%%%%%%%%%%%%%%%%%%%%%%%%%%

%%%%%%%%%%%%%%%%%%%%%%%%%%%%%%%%%%%%%%%%%%%%%%%%%%%%%%%%%%%%%%%%%%%%%%%%%%%%%%%
% \begin{frame}{The hidden cost of benchmarking}
%
%     Different teams often rebuild the same engineering stack:\\[.5em]
%     \begin{itemize}
%         \item Install and environment creation.
%         \item Dependency management and data preparation.
%         \item Run orchestration, evaluation and plotting.
%     \end{itemize}
%
%     \vskip1.5em
%     \strongpoint{Repeated tooling slows down science and creates avoidable burden.}
%
% \end{frame}
%%%%%%%%%%%%%%%%%%%%%%%%%%%%%%%%%%%%%%%%%%%%%%%%%%%%%%%%%%%%%%%%%%%%%%%%%%%%%%%




%%%%%%%%%%%%%%%%%%%%%%%%%%%%%%%%%%%%%%%%%%%%%%%%%%%%%%%%%%%%%%%%%%%%%%%%%%%%%%%
% \begin{frame}{What \Benchopt{} changes}
%
%     \Benchopt{} factorizes the common benchmark tooling:\\[.5em]
%     \begin{itemize}
%         \item Standard interfaces: objective, dataset, solver.
%         \item Reusable workflow: install, run, test, compare, share.
%         \item Run anywhere: explicit dependencies and controlled environments.
%         \item Reliable reruns: controlled randomness and reproducible protocols.
%         \item New in recent releases: merge runs and publish to GitHub/Hugging Face.
%     \end{itemize}
%
%     \pause
%     \vskip1.5em
%     \strongpoint{Benchopt does not replace benchmarks,\\it removes duplicated plumbing.}
%
% \end{frame}
%%%%%%%%%%%%%%%%%%%%%%%%%%%%%%%%%%%%%%%%%%%%%%%%%%%%%%%%%%%%%%%%%%%%%%%%%%%%%%%


%%%%%%%%%%%%%%%%%%%%%%%%%%%%%%%%%%%%%%%%%%%%%%%%%%%%%%%%%%%%%%%%%%%%%%%%%%%%%%%
% \begin{frame}{Benchopt scope}
%
%     \Benchopt{} turns benchmark code into reusable research infrastructure.
%     \vskip1.5em
%
%     {\centering
%     \vskip2em
%     {\Large\Benchopt{} produces \textbf{open}, \textbf{reproducible}, \textbf{extendable} benchmarks.}\\
%     }
%     \vskip1em
%     \begin{itemize}
%         \item Comparable and fairer evaluations,
%         \item Faster iterations on methods and datasets,
%         \item Lower engineering overhead across projects.
%     \end{itemize}
%
%     \vskip1.2em
%     {\centering And keeps them \textbf{maintainable} through explicit interfaces\\and tests.\\}
%
% \end{frame}
%%%%%%%%%%%%%%%%%%%%%%%%%%%%%%%%%%%%%%%%%%%%%%%%%%%%%%%%%%%%%%%%%%%%%%%%%%%%%%%



%%%%%%%%%%%%%%%%%%%%%%%%%%%%%%%%%%%%%%%%%%%%%%%%%%%%%%%%%%%%%%%%%%%%%%%%%%%%%%%
% \begin{frame}{Extra tools for easier benchmarking}
%
%     \begin{itemize}
%         \item \textbf{Integrate broadly:} use implementations from Python, R, Julia, or binaries.
%         \item \textbf{Scale experiments:} run in parallel locally or on distributed backends.
%         \item \textbf{Save time:} cache results to avoid recomputing unchanged runs.
%         \item \textbf{Trust comparisons:} control randomness with seeds and stable protocols.
%         \item \textbf{Share outcomes:} merge and publish results from multiple runs, easy visualization.
%     \end{itemize}
%
%     \vskip0.7em
%     \centering
%     \includegraphics[width=0.8\linewidth]{tweet_rahimi.png}
%
% \end{frame}
%%%%%%%%%%%%%%%%%%%%%%%%%%%%%%%%%%%%%%%%%%%%%%%%%%%%%%%%%%%%%%%%%%%%%%%%%%%%%%%



%%%%%%%%%%%%%%%%%%%%%%%%%%%%%%%%%%%%%%%%%%%%%%%%%%%%%%%%%%%%%%%%%%%%%%%%%%%%%%%
% \begin{frame}{How does \Benchopt{} do it?}
%
% \Benchopt{} provides a framework to organize and run benchmarks:
% \begin{itemize}
%     \item Each benchmark is a repository
%     \item \Benchopt{} is a \texttt{Python} CLI to install, run, test and share them
% \end{itemize}
%
% \vskip1em
% Examples of existing benchmarks:
% \vskip.5em
% \begin{minipage}{0.5\linewidth}
%     {\small
%     \begin{itemize}
%         \item \href{https://github.com/benchopt/benchmark_resnet_classif}{\bf Image Classification (resnet)}
%         \item \href{https://github.com/benchopt/benchmark_logreg_l2}{\bf Logistic regression}
%         \item \href{https://github.com/benchopt/benchmark_lasso}{\bf Lasso}
%         \item \href{https://github.com/benchopt/benchmark_linear_ica}{\bf ICA}
%     \end{itemize}
%     }
% \end{minipage}
% \begin{minipage}{0.48\linewidth}
%     {\small
%     \begin{itemize}
%         \item \href{https://github.com/benchopt/benchmark_tv_1d}{\bf Total Variation}
%         \item \href{https://github.com/benchopt/benchmark_bilevel}{\bf Bilevel Optimization}
%         \item \href{https://github.com/benchopt/benchmark_bci}{\bf Brain Computer Interface}
%         \item \href{https://github.com/benchopt/?q=benchmark_&type=all&language=&sort=name}{\bf \dots}
%     \end{itemize}
%     }
% \end{minipage}
% \vskip2em
%
% \centering
% Start yours with \url{https://github.com/benchopt/template_benchmark} or \url{https://github.com/benchopt/template_benchmark_ml}!
% \end{frame}
%%%%%%%%%%%%%%%%%%%%%%%%%%%%%%%%%%%%%%%%%%%%%%%%%%%%%%%%%%%%%%%%%%%%%%%%%%%%%%%


% -----------------------------------------------------------------------------
% TIMING NOTES (1h50 hands-on)
% - ~15 min : Example 1 — Denoising: run existing + add a solver
% - ~40 min : Participants add their solver, run benchopt
% - ~15 min : Example 2 — Blind deblurring: objective skeleton walkthrough
% - ~35 min : Participants write objective.py, run both solvers
% -  ~5 min : Wrap-up, share solutions, Q&A
%
% IF SHORT ON TIME
% - Keep: install/run workflow + 3 solver steps + 2 objective steps
% - Skip: Components/Flow/Structure details + long benchmark example list
% -----------------------------------------------------------------------------


%%%%%%%%%%%%%%%%%%%%%%%%%%%%%%%%%%%%%%%%%%%%%%%%%%%%%%%%%%%%%%%%%%%%%%%%%%%%%%%
\begin{frame}{Tutorial}

    \centering
    \huge Sorry I failed to release for the tutorial....\\[2em]

    \pause

    But you will get to ride on the edge!\\[1em]

    \large \texttt{pip install -U git+https://github.com/benchopt/benchopt}

    \vskip2em
    You will just need an env with deep-inverse, benchopt and optinonally ptwt for the wavelet baseline.\\[1em]

\end{frame}
%%%%%%%%%%%%%%%%%%%%%%%%%%%%%%%%%%%%%%%%%%%%%%%%%%%%%%%%%%%%%%%%%%%%%%%%%%%%%%%



%%%%%%%%%%%%%%%%%%%%%%%%%%%%%%%%%%%%%%%%%%%%%%%%%%%%%%%%%%%%%%%%%%%%%%%%%%%%%%%
\begin{frame}{Tutorial}

    \centering
    \huge Running a benchmark!

\end{frame}
%%%%%%%%%%%%%%%%%%%%%%%%%%%%%%%%%%%%%%%%%%%%%%%%%%%%%%%%%%%%%%%%%%%%%%%%%%%%%%%


%%%%%%%%%%%%%%%%%%%%%%%%%%%%%%%%%%%%%%%%%%%%%%%%%%%%%%%%%%%%%%%%%%%%%%%%%%%%%%%
\begin{frame}[fragile]{Install and run the benchmark}

    \textbf{Repository used in the tutorial:}\\[.3em]
    \url{\TutorialRepoURL}

    \vskip1em
    \textbf{Minimal workflow:}

\begin{lstlisting}[language=bash]
git clone https://github.com/tomMoral/26-05_benchopt_tuto_icms tuto_benchopt
cd tuto_benchopt/01-denoising
pip install benchopt
benchopt install .
benchopt run .
\end{lstlisting}

    \vskip0.7em
    \strongpoint{Once \Benchopt{} is installed, two commands are enough: install the benchmark, then run it.}

\end{frame}
%%%%%%%%%%%%%%%%%%%%%%%%%%%%%%%%%%%%%%%%%%%%%%%%%%%%%%%%%%%%%%%%%%%%%%%%%%%%%%%


%%%%%%%%%%%%%%%%%%%%%%%%%%%%%%%%%%%%%%%%%%%%%%%%%%%%%%%%%%%%%%%%%%%%%%%%%%%%%%%
\begin{frame}{What happens during install and run?}

    For this benchmark, \Benchopt{} will automatically:\\[.5em]
    \begin{itemize}
        \item Load the objective, dataset and available solvers.
        \item Install the benchmark requirements declared in the files.
        \item Run the available solver on the \texttt{CBSD} dataset.
        \item Collect metrics such as \texttt{psnr} and \texttt{ssim}.
    \end{itemize}

    \vskip1.2em
    \textbf{In this repo:}
    \begin{itemize}
        \item Dataset: \texttt{datasets/cbsd.py} — CBSD images, parametrized by $\sigma$
        \item Objective: \texttt{objective.py} — adds noise, computes PSNR
        \item Solvers: \texttt{solvers/wavelet.py}, \texttt{solvers/dncnn.py}
        \item Plot: \texttt{plots/denoised\_images.py} -- display results
    \end{itemize}

\end{frame}
%%%%%%%%%%%%%%%%%%%%%%%%%%%%%%%%%%%%%%%%%%%%%%%%%%%%%%%%%%%%%%%%%%%%%%%%%%%%%%%


%%%%%%%%%%%%%%%%%%%%%%%%%%%%%%%%%%%%%%%%%%%%%%%%%%%%%%%%%%%%%%%%%%%%%%%%%%%%%%%
\begin{frame}{Components of a benchmark}

    \centering
    \vspace*{-5mm}
    \textbf{3 components}: Objective, Dataset, Solver

    \vskip3em
    \includegraphics[width=\linewidth]{benchopt_schema.pdf}

\end{frame}
%%%%%%%%%%%%%%%%%%%%%%%%%%%%%%%%%%%%%%%%%%%%%%%%%%%%%%%%%%%%%%%%%%%%%%%%%%%%%%%


%%%%%%%%%%%%%%%%%%%%%%%%%%%%%%%%%%%%%%%%%%%%%%%%%%%%%%%%%%%%%%%%%%%%%%%%%%%%%%%
\begin{frame}{Tutorial — Example 1}

    \centering
    \huge Adding a denoiser solver!

    \vskip1em
    {\large\textit{Benchmark:} \texttt{denoising}}

\end{frame}
%%%%%%%%%%%%%%%%%%%%%%%%%%%%%%%%%%%%%%%%%%%%%%%%%%%%%%%%%%%%%%%%%%%%%%%%%%%%%%%


%%%%%%%%%%%%%%%%%%%%%%%%%%%%%%%%%%%%%%%%%%%%%%%%%%%%%%%%%%%%%%%%%%%%%%%%%%%%%%%
\begin{frame}{Structure of a benchmark}

    \begin{minipage}{0.45\linewidth}
        \dirtree{%
        .1 benchmark/.
        .2 objective.py.
        .2 datasets/.
        .3 dataset1.py.
        .3 dataset2.py.
        .2 solvers/.
        .3 solver1.py.
        .3 solver2.py.
    }
    \end{minipage}
    \begin{minipage}{0.45\linewidth}

    \textbf{Modular \& extendable}
    \vspace*{5mm}

    New solver? add a file

    New dataset? add a file

    New metric? modify objective
    \end{minipage}

\end{frame}
%%%%%%%%%%%%%%%%%%%%%%%%%%%%%%%%%%%%%%%%%%%%%%%%%%%%%%%%%%%%%%%%%%%%%%%%%%%%%%%


%%%%%%%%%%%%%%%%%%%%%%%%%%%%%%%%%%%%%%%%%%%%%%%%%%%%%%%%%%%%%%%%%%%%%%%%%%%%%%%
\frame[t]{
    \frametitle{A modular framework to create benchmarks}

    \centering
    \vskip.5em
    \textbf{3 components}: Objective, Dataset, Solver

    \vskip2em
    \includegraphics[width=\linewidth]{benchopt_schema_dependency.png}

    \pause
    \strongpoint{Benchopt defines the interface between components.}

}
%%%%%%%%%%%%%%%%%%%%%%%%%%%%%%%%%%%%%%%%%%%%%%%%%%%%%%%%%%%%%%%%%%%%%%%%%%%%%%%


%%%%%%%%%%%%%%%%%%%%%%%%%%%%%%%%%%%%%%%%%%%%%%%%%%%%%%%%%%%%%%%%%%%%%%%%%%%%%%%
\frame[t]{
    \frametitle{Explicit preprocessing}

    \inputpython{benchopt_dataset.py}{1}{19}

    \pause
    \strongpoint{Reproducible benchmark by design!}

    \vskip1em
    \textbf{Goal:} if you use the same setup, you don't need to re-run the baseline methods!

}
%%%%%%%%%%%%%%%%%%%%%%%%%%%%%%%%%%%%%%%%%%%%%%%%%%%%%%%%%%%%%%%%%%%%%%%%%%%%%%%


%%%%%%%%%%%%%%%%%%%%%%%%%%%%%%%%%%%%%%%%%%%%%%%%%%%%%%%%%%%%%%%%%%%%%%%%%%%%%%%
\frame[t]{
    \frametitle{Explicit evaluation protocol and requirements}

    \inputpython{benchopt_objective.py}{1}{29}

}
%%%%%%%%%%%%%%%%%%%%%%%%%%%%%%%%%%%%%%%%%%%%%%%%%%%%%%%%%%%%%%%%%%%%%%%%%%%%%%%


%%%%%%%%%%%%%%%%%%%%%%%%%%%%%%%%%%%%%%%%%%%%%%%%%%%%%%%%%%%%%%%%%%%%%%%%%%%%%%%
\frame[t]{
    \frametitle{Explicit parameters}

    \inputpython{benchopt_solver.py}{1}{26}

}
%%%%%%%%%%%%%%%%%%%%%%%%%%%%%%%%%%%%%%%%%%%%%%%%%%%%%%%%%%%%%%%%%%%%%%%%%%%%%%%




%%%%%%%%%%%%%%%%%%%%%%%%%%%%%%%%%%%%%%%%%%%%%%%%%%%%%%%%%%%%%%%%%%%%%%%%%%%%%%%
\begin{frame}[fragile]{Step 1: add a new solver file}

    \textbf{Goal:} add a new deep denoiser without changing the benchmark protocol.

    \vskip1em
    \begin{itemize}
        \item Open \texttt{solvers/skeleton.py} — your starting point
        \item Rename the solver (\texttt{name} attribute) to \NewSolver{}
        \item Instantiate your denoiser in \texttt{run()}, e.g.:
\begin{lstlisting}[language=Python]
dinv.models.DnCNN(pretrained='download', device=self.device)
\end{lstlisting}
        \item Keep \texttt{sampling\_strategy = 'run\_once'} — no loop needed
        \item Save the output in \texttt{self.x\_hat}
    \end{itemize}

    \strongpoint{Same interface, same objective, same dataset: only the model changes.}

\end{frame}
%%%%%%%%%%%%%%%%%%%%%%%%%%%%%%%%%%%%%%%%%%%%%%%%%%%%%%%%%%%%%%%%%%%%%%%%%%%%%%%


%%%%%%%%%%%%%%%%%%%%%%%%%%%%%%%%%%%%%%%%%%%%%%%%%%%%%%%%%%%%%%%%%%%%%%%%%%%%%%%
\begin{frame}[fragile]{Step 2: plug in your denoiser}

    All deepinv denoisers share the same calling convention:\\[.5em]
    \begin{center}
    \texttt{x\_hat = denoiser(y, sigma)}
    \end{center}

    \vskip0.8em
    Some models to try (all available with pretrained weights):

    \vskip0.5em
    \begin{tabular}{lll}
        \textbf{Model} & \textbf{Type} & \textbf{Key property} \\
        \hline
        \texttt{DnCNN}  & deep, fixed $\sigma$ & fast, good baseline \\
        \texttt{SCUNet} & deep, adaptive $\sigma$ & robust across levels \\
        \texttt{SwinIR} & transformer & state-of-the-art \\
    \end{tabular}

    \vskip1em
    \strongpoint{Explore \href{https://deepinv.github.io/deepinv/auto_examples/models/demo_denoiser_tour.html}{deepinv denoiser tour} for the full list.}

\end{frame}
%%%%%%%%%%%%%%%%%%%%%%%%%%%%%%%%%%%%%%%%%%%%%%%%%%%%%%%%%%%%%%%%%%%%%%%%%%%%%%%


%%%%%%%%%%%%%%%%%%%%%%%%%%%%%%%%%%%%%%%%%%%%%%%%%%%%%%%%%%%%%%%%%%%%%%%%%%%%%%%
\begin{frame}[fragile]{Step 3: validate and compare}

    \Benchopt{} automatically discovers the new solver and compares it.

    \vskip1em
\begin{lstlisting}[language=bash]
benchopt test . -k mysolver --skip-env
benchopt run . -s MySolver -s DRUNet -s WaveletDict
\end{lstlisting}

    \vskip0.7em
    What to check:
    \begin{itemize}
        \item The solver file is picked up automatically.
        \item \texttt{get\_result()} returns \texttt{\{"x\_hat": ...\}}.
        \item \texttt{psnr} is reported alongside the existing solvers.
        \item No change needed in the objective or dataset.
    \end{itemize}

    \vskip0.7em
    \strongpoint{Adding a solver means writing one compatible file — nothing else changes.}

\end{frame}
%%%%%%%%%%%%%%%%%%%%%%%%%%%%%%%%%%%%%%%%%%%%%%%%%%%%%%%%%%%%%%%%%%%%%%%%%%%%%%%



%%%%%%%%%%%%%%%%%%%%%%%%%%%%%%%%%%%%%%%%%%%%%%%%%%%%%%%%%%%%%%%%%%%%%%%%%%%%%%%
\begin{frame}{Tutorial — Example 2}

    \centering
    \huge Writing the objective\\ for blind deblurring!

    \vskip1em
    {\large\textit{Benchmark:} \texttt{blind\_deblur}}

\end{frame}
%%%%%%%%%%%%%%%%%%%%%%%%%%%%%%%%%%%%%%%%%%%%%%%%%%%%%%%%%%%%%%%%%%%%%%%%%%%%%%%


%%%%%%%%%%%%%%%%%%%%%%%%%%%%%%%%%%%%%%%%%%%%%%%%%%%%%%%%%%%%%%%%%%%%%%%%%%%%%%%
\begin{frame}{Step 1: understand what is given}

    \textbf{Goal:} implement \texttt{objective.py} for a blind deblurring benchmark.
    The solvers already exist — only the objective is missing.

    \vskip1em
    \textbf{Given for free} (non-obvious, provided in the skeleton):
    \begin{itemize}
        \item \texttt{set\_data()} — stores $x$, $y$, physics, \texttt{kernel\_size}
        \item \texttt{get\_one\_result()} — delta kernel baseline
    \end{itemize}

    \vskip0.7em
    \textbf{Yours to implement:}
    \begin{itemize}
        \item \texttt{get\_objective()} — what to expose to solvers
        \item \texttt{evaluate\_result()} — dual metric (PSNR + kernel MSE)
        \item \texttt{save\_final\_results()} — persist both \texttt{x\_hat} and \texttt{k\_hat}
    \end{itemize}

    \strongpoint{The \texttt{while callback()} pattern is illustrated by \texttt{dip\_selfdeblur.py}.}

\end{frame}
%%%%%%%%%%%%%%%%%%%%%%%%%%%%%%%%%%%%%%%%%%%%%%%%%%%%%%%%%%%%%%%%%%%%%%%%%%%%%%%


%%%%%%%%%%%%%%%%%%%%%%%%%%%%%%%%%%%%%%%%%%%%%%%%%%%%%%%%%%%%%%%%%%%%%%%%%%%%%%%
\begin{frame}[fragile]{Step 2: implement the objective methods}

    \textbf{\texttt{get\_objective}:} expose \texttt{y} and \texttt{kernel\_size} — do NOT expose the true kernel.

    \vskip0.7em
    \textbf{\texttt{evaluate\_result}:} dual metric
    \begin{itemize}
        \item Image quality: $\mathrm{PSNR}(\hat{x}, x_{\mathrm{true}})$
        \item Kernel accuracy: normalise both kernels, then compute MSE
        $$\mathrm{kernel\_mse} = \left\| \frac{\hat{k}}{\|\hat{k}\|_1} - \frac{k}{\|k\|_1} \right\|_2^2$$
    \end{itemize}

    \vskip0.5em
    \textbf{\texttt{save\_final\_results}:} persist both estimates to CPU
\begin{lstlisting}[language=Python]
return dict(x_hat=x_hat.cpu(), k_hat=k_hat.cpu())
\end{lstlisting}

    \strongpoint{The stopping criterion monitors \texttt{psnr} — make sure it is in the returned dict.}

\end{frame}
%%%%%%%%%%%%%%%%%%%%%%%%%%%%%%%%%%%%%%%%%%%%%%%%%%%%%%%%%%%%%%%%%%%%%%%%%%%%%%%


%%%%%%%%%%%%%%%%%%%%%%%%%%%%%%%%%%%%%%%%%%%%%%%%%%%%%%%%%%%%%%%%%%%%%%%%%%%%%%%
\begin{frame}[fragile]{Step 3: validate the objective}

    \Benchopt{} runs both solvers and checks your objective.

    \vskip1em
\begin{lstlisting}[language=bash]
benchopt test . --skip-env
benchopt run . -s Bilevel-PGD -s DIP-SelfDeblur
\end{lstlisting}

    \vskip0.7em
    What to check:
    \begin{itemize}
        \item \texttt{psnr} and \texttt{kernel\_mse} are reported for both solvers.
        \item The convergence curve improves over the delta-kernel baseline.
        \item The \textbf{Kernel} plot shows the estimated blur kernel per solver.
        \item Both solvers complete within the time budget.
    \end{itemize}

    \vskip0.7em
    \strongpoint{Writing an objective means defining the protocol once —\\both solvers run against it automatically.}

\end{frame}
%%%%%%%%%%%%%%%%%%%%%%%%%%%%%%%%%%%%%%%%%%%%%%%%%%%%%%%%%%%%%%%%%%%%%%%%%%%%%%%


%%%%%%%%%%%%%%%%%%%%%%%%%%%%%%%%%%%%%%%%%%%%%%%%%%%%%%%%%%%%%%%%%%%%%%%%%%%%%%%
\begin{frame}{\Benchopt{} Roadmap}
    \begin{itemize}\itemsep1em
        \item \textbf{Simpler solver evaluation --} make custom solver integration more direct, for instance with \texttt{eval\_solver(run\_func)}
        \item \textbf{Better visualization tools -- } explore results more easily, in particular for image tasks
        \item \textbf{More reference benchmarks --} grow the benchmark collection with community contributions
    \end{itemize}


    \strongpoint{\Benchopt{} is by no means finished or complete!}

    \vskip2em
    It is useful to a growing group of people but we need to gather more feedbacks to fit the needs of the community.\\[2em]
\end{frame}
%%%%%%%%%%%%%%%%%%%%%%%%%%%%%%%%%%%%%%%%%%%%%%%%%%%%%%%%%%%%%%%%%%%%%%%%%%%%%%%


%%%%%%%%%%%%%%%%%%%%%%%%%%%%%%%%%%%%%%%%%%%%%%%%%%%%%%%%%%%%%%%%%%%%%%%%%%%%%%%
\begin{frame}{Benchopt links and resources}

    \begin{columns}[T]
        \column{.5\textwidth}
        \vskip1em
        \begin{itemize}\itemsep1em
            \item Documentation:\\ \href{https://benchopt.github.io}{benchopt.github.io}
            \item Main repository: \href{https://github.com/benchopt/benchopt}{github.com/benchopt/benchopt}
            \item Benchmark templates: \href{https://github.com/benchopt/template_benchmark}{github.com/benchopt/template\_benchmark}\\
            \href{https://github.com/benchopt/template_benchmark_ml}{github.com/benchopt/template\_benchmark\_ml}
        \end{itemize}

        \column{.36\textwidth}
        \centering
        \includegraphics[width=.8\textwidth]{qr_benchopt_repos.png}\\[.5em]
    \end{columns}
    \vskip1.2em
    \strongpoint{If this tutorial was useful, please star the \texttt{benchopt} repo and share it.}

\end{frame}
%%%%%%%%%%%%%%%%%%%%%%%%%%%%%%%%%%%%%%%%%%%%%%%%%%%%%%%%%%%%%%%%%%%%%%%%%%%%%%%

\end{document}
